\subsection{A finite identity converting nesting into visits}
\label{sec:finite-visits}
The nesting estimate counts polygons.  We now insert a visited vertex
using an identity in a finite domain.  No limiting path law is needed. The same finite recursion appears in
\cite[Proposition~8.1]{PolynomialVacuumCompanion}; we include its proof
because positivity on filled sets is needed after deleting a chord.

A finite vertex set $D$ is \emph{filled} if every lattice vertex strictly
inside a simple polygon in its induced graph also belongs to $D$.
Its boundary ports are midpoints of edges with one endpoint in $D$ and
the other outside.  Sum excursions with both orders on their distinct
endpoints.  For $v\in D$, define
\[
 T(v,D)=\sum_{\substack{\gamma\text{ excursion in }D\\v\in\gamma}}
       x^{\ell(\gamma)}e^{isW(\gamma)},\qquad s=3/8,
\]
where $W$ is the signed sum of the turns, including those at the first
and last visited vertices.  Reversal conjugates each term, so $T(v,D)$
is real.  For a general filled set its individual terms need not have
positive real part.

For a face center $f$, let $Z(D;f)$ sum collections of pairwise
vertex-disjoint simple polygons in $D$ enclosing $f$, with weight
$2x^{|\Gamma|}$ per polygon and empty weight one.  For a polygon $\Gamma$,
write $\Gamma^\circ$ for its strict interior vertex set.  If $v\in\Gamma$,
let $\tau_v(\Gamma)=\pm\pi/3$ be its turn at $v$ in counterclockwise
traversal.

\begin{proposition}[Visit recursion and nesting comparison]
\label{prop:visits}
For every finite filled $D$ and $v\in D$,
\begin{align}
T(v,D)=3
 &+\sum_{\substack{\Gamma\subset D\\v\in\Gamma}}
       2x^{|\Gamma|-1}\big[-\cos\big(s(2\pi-\tau_v(\Gamma))\big)\big]
 \notag\\
 &+\sum_{\substack{\Gamma\subset D\\v\in\Gamma^\circ}}
       2x^{|\Gamma|}T(v,\Gamma^\circ).                 \label{c2:eq-G2}
\end{align}
Every coefficient on the right is positive.  In particular, $T(v,D)\ge3$
and $T(v,D)$ is monotone under inclusion of finite filled sets containing
$v$.  There are absolute positive constants such that
\begin{equation}
 T(v,D)\asymp\sum_{f\text{ incident to }v}Z(D;f).       \label{c2:eq-G3}
\end{equation}
If $D$ is the vertex set inside a bounded convex polygon cut out by
admissible walls, and $A_D(v)$ denotes the corresponding positive,
ordered excursion weight through $v$, then
\begin{equation}
 \cos(3\pi/8)A_D(v)\le T(v,D)\le A_D(v).              \label{eq:positive-visits}
\end{equation}
\end{proposition}

\begin{proof}
On arriving at a fresh vertex along one edge, take either nonbacktracking
continuation with weight $xe^{is\pi/3}$ or $xe^{-is\pi/3}$.  These weights
sum to one.  Stop at a port exit or before making a choice at the first
repeated vertex.  The exploration tree is finite; replacing a node by
its children preserves total weight.  Consequently all terminal weights
sum to one.  A prefix stopped upon arrival at a vertex has not yet taken
that vertex's weight or turn.

Start at a boundary port of $D$.  A first repeat consists of a simple
stem and a simple attached polygon.  The stem reaches the polygon from
its exterior.  Otherwise the outside neighbor of its initial port would
lie strictly inside that polygon: the stem could not leave the interior
without previously hitting the polygon.  This contradicts filling.
Reversing the attached loop changes its added turn from $4\pi/3$ to
$-4\pi/3$.  The two phases sum to
$2\cos(4\pi s/3)=0$.  This pairing preserves the visited vertices,
so it still cancels when we retain only histories that have visited $v$.
Continuing any first-arrival prefix at $v$ has total weight one.
Thus the sum of such prefixes, over all initial ports, equals $T(v,D)$.

Next start at $v$ along each of its three edges with weight one and no
turn or vertex weight at $v$.  Stop at an exit, a return to $v$, or the
first repeat of another vertex.  These three experiments have total
weight three.  Exit terms are reversals of the first-arrival prefixes,
so their sum is $\overline{T(v,D)}=T(v,D)$.

A return to $v$ traverses a simple polygon $\Gamma$ through it.  The
counterclockwise turn omits the turn at $v$ and equals
$2\pi-\tau_v(\Gamma)$; the reverse orientation has opposite turn.
Their total terminal contribution is
\[
 2x^{|\Gamma|-1}\cos\big(s(2\pi-\tau_v(\Gamma))\big).
\]
The missing factor of $x$ is exactly the unweighted starting vertex.

For a repeat at another vertex, the attached polygon avoids $v$.
If $v$ is outside, the same exterior cancellation applies.
If $v$ is inside, its stem reaches the attachment along the interior
third edge.  The turn there in counterclockwise traversal is $-\pi/3$,
whereas arriving from the stem gives $+\pi/3$.  The two possible added
turns are therefore $8\pi/3$ and $-8\pi/3$, both with phase $-1$.
The combined factor is $-2x^{|\Gamma|}$.

For a fixed polygon enclosing $v$, these stems are exactly the simple
paths from $v$ exiting $\Gamma^\circ$, with no weight or turn at $v$.
Filling gives $\Gamma^\circ\subset D$, and planarity implies that an
edge leaving the strict interior ends on $\Gamma$.  Completing the
last half-edge therefore reaches the attachment without adding another
weighted vertex.  The stem sum is $T(v,\Gamma^\circ)$, by the prefix
reversal just proved in that smaller filled domain.  Summing all terminal
contributions and moving the non-exit terms to the other side proves
\eqref{c2:eq-G2}.

To prove positivity and the comparison, there are only two local cases.
A through-polygon encloses one incident face center when its turn at
\(v\) is \(+\pi/3\), and two when that turn is \(-\pi/3\). The omitted-turn
angles are respectively \(5\pi/3\) and \(7\pi/3\). In these cases set
\[
 c_1=-\cos(5\pi/8)>0,\qquad c_2=-\cos(7\pi/8)>0.
\]
The recursion terminates, since taking strict interior reduces the number
of vertices.  Unroll it into an outer-to-inner chain of disjoint polygons
strictly enclosing $v$.  Each chain carries the product of factors
$2x^{|\Gamma|}$.  In its remaining interior $D'$, the last factor is
\[
 3+\frac2x\sum_{\substack{Q\subset D'\\v\in Q}}
             c_{m_Q}x^{|Q|},\qquad m_Q\in\{1,2\}.
\]
This positive expansion gives both $T(v,D)\ge3$ and monotonicity: inclusion
only adds eligible chains and through-polygons.

The same chains enumerate the three incident-face partition functions.
A polygon strictly enclosing $v$ surrounds all three incident face
centers.  A polygon through $v$ surrounds precisely $m_Q$ of them.
Conversely, a polygon surrounding one incident center has $v$ inside
or on its boundary, since the open face approaches $v$ and does not meet
the polygon.  At most one polygon of a disjoint collection can pass
through $v$, and that polygon is innermost.  Hence the final factor for
the incident-face sum is
\[
 3+2\sum_{\substack{Q\subset D'\\v\in Q}}m_Qx^{|Q|}.
\]
The ratios $c_m/(xm)$, for $m=1,2$, are fixed positive numbers.
Termwise comparison proves \eqref{c2:eq-G3}.

Finally, for an admissible convex domain the turn rule proved with
\eqref{c2:eq-1} gives $-\pi\le W\le\pi$.  Every real part is therefore
between $\cos(3\pi/8)$ and one times its positive path weight.
This proves \eqref{eq:positive-visits}.
\end{proof}

Proposition~\ref{prop:visits} has completed the conversion from polygon
nesting to bulk-visit weight.  It has not chosen the path's terminal side.
We first use it for the triangle law, where both endpoints are summed;
the subsequent sewing arguments select arches and then bridges.
